\section{Principal data, admissible windows, and exact character extraction}
\label{sec:sp4-principal-window}

We separate three assertions: a source-dependent automorphic realization,
an abstract section-selection argument, and an exact arithmetic consumer.
Only the algebraic power-section and finite-observation inputs below are
newly kernel checked. None is an independent cancellation estimate.

\paragraph{The principal Hecke character is realized.}
Let \(K\supset\mathbb Q(\mu_8)\). For the simply connected fourfold
\(\operatorname{Sp}_4\) cover, choose the diagonal bisector
\(D(e_i,e_j)=-\delta_{ij}\) and the trivial Brylinski--Deligne
root-lift parameter. These are the Siegel-Levi parameters \(p=-1,q=0\)
in \cite[Section 6]{GaoInverseEis}.
The torus cocycle is
\[
 c(t,u)=\prod_{i=1}^2(t_i,u_i)_4^{-1}.
\]
Here \(Y_{Q,4}=2Y\), so its centre lies over square coordinates.
Bimultiplicativity gives
\((a^2,b^2)_4=(a,b)_4^4=1\): the chosen section is a homomorphism
on this central torus. Abstractly, if a section has bimultiplicative
defect of exponent dividing \(n\), its restriction along the
\(m\)-power map is multiplicative whenever \(n\mid m^2\).
This is the maintained theorem
\texttt{Metaplectic.Sp4FourInverse.power\_section\_exists}; its
bimultiplicative power calculation is
\texttt{Metaplectic.Sp4FourInverse.power\_pair}.
It does not itself construct a covering group or an adelic quotient.

On the global central torus define
\[
 \chi_0(\zeta,t_1,t_2)
   =\epsilon(\zeta)\lvert t_1/t_2\rvert_{\mathbb A}^{1/4}.
\]
This is genuine and continuous. The product formula makes it trivial
on the rational torus in the centre, so it is automorphic in the
convention of \cite[Section 2.6]{GaoInverseEis}.
The root lifts, with the chosen trivial root-lift parameter, give
\[
 \chi_0(h_{\alpha_1}(a^2))=|a|_{\mathbb A},\qquad
 \chi_0(h_{\alpha_2}(a^4))=|a|_{\mathbb A}^{-1}.
\]
Indeed the two torus coordinates are \((a^2,a^{-2})\) and
\((1,a^4)\), respectively. The datum is Siegel-Levi exceptional and
unitary on the scalar centre. It is unramified outside finitely many
places, which we include in \(S\). The source distinguishedness and
factorization theorems therefore yield
\begin{equation}\label{eq:sp4-principal-inverse}
 E_\psi(s)=\frac{B(s)}{\zeta_K^S(4s+2)}.
\end{equation}
Thus \(\eta=1\) is available; it is not necessary to assume a fourth
root of a desired Dirichlet character on a fixed field.
This application of Gao is an ordinary source-dependent argument,
not a Lean formalization of his automorphic theorems.

\paragraph{A finite-dimensional observation lemma.}
Let \(E\) be a topological complex vector space, \(E_{\rm fin}\)
a dense linear subspace, and \(F:E\to V\) a continuous linear map
to a finite-dimensional complex normed space. Then
\[
 F(E_{\rm fin})=F(E).
\]
For \(R=F(E_{\rm fin})\), finite dimensionality makes \(R\) closed.
Its closed preimage contains the dense subspace, hence equals \(E\).
This full proof is maintained as
\texttt{Metaplectic.DenseObservation.exists\_same\_observation}.
The Lean theorem allows any complete nontrivially normed scalar field;
it imposes no norm or completeness on \(E\).

\paragraph{Finite-point nonvanishing with \(K\)-finite data.}
The following analytic interface is explicit. Work in a fixed smooth
compact-picture space \(E\); its \(K\)-finite vectors form the dense
linear subspace \(E_{\rm fin}\). Suppose the local Jacquet family is
a meromorphic family of continuous linear functionals on \(E\), with
locally finite pole order in that topology. Near \(s_0\), write
\[
 W(s)=\sum_{\ell=-r}^{\infty}\lambda_\ell(s-s_0)^\ell,
 \qquad \lambda_\ell\in E',\quad r\ge0.
\]
Also suppose the usual open-Bruhat-cell construction supplies a
holomorphic smooth section \(h(s)\) with \(W(s)h(s)=1\).
One takes a Levi vector on which its Whittaker functional is \(1\),
and a compactly supported smooth function \(\phi\) on the integration
unipotent satisfying \(\int\phi(u)\psi(u)^{-1}\,du=1\).
These are representation-theoretic hypotheses, not conclusions of
the abstract Lean lemma. Gao's local integral in Section 4.2 is the
relevant interface; using its continuous smooth meromorphic realization
is an additional analytic input here.

Write \(h(s)=\sum_{j\ge0}h_j(s-s_0)^j\). Apply the observation lemma
to \(F=(\lambda_{-r},\ldots,\lambda_0)\) and each \(h_j\),
\(0\le j\le r\). Obtain \(k_j\in E_{\rm fin}\) with the same
observations, and put
\[
 P(s)=\sum_{j=0}^r k_j(s-s_0)^j.
\]
For every \(-r\le a\le0\), the coefficient of \((s-s_0)^a\)
in \(W(s)P(s)\) is
\(\sum_{0\le j\le a+r}\lambda_{a-j}(k_j)\).
It equals the corresponding coefficient for \(W(s)h(s)=1\).
Consequently \(W(s)P(s)\) is holomorphic at \(s_0\) and equals
\(1\) there. The section \(P\) lies in the span of finitely many
\(K\)-types and is a polynomial combination of admissible flat
sections. This avoids inserting a smooth non-\(K\)-finite test directly
into Gao's \(K\)-finite Eisenstein-series statement.

For distinct finitely many points \(s_i\), with pole orders at most
\(r_i\), make these finite observations at all points. Vector-valued
Hermite interpolation gives one polynomial section of degree at most
\(\sum_i(r_i+1)-1\) with the prescribed jets. The same calculation
gives value \(1\) at each point. Taking the finite product over bad
places makes \(B\) holomorphic with \(B(s_i)=1\), under the stated
local analytic hypotheses. No bound on the selected \(K\)-types,
coefficients or Sobolev norms as height or conductor grows is asserted.

\paragraph{Extract a chosen Dirichlet character without discarding zeros.}
Let \(H\) be a finite subgroup of Dirichlet characters containing
the four characters modulo \(8\), and let \(\chi\in H\).
Use a common modulus and let \(K\) be the abelian field whose
character group is \(H\). Then \(K\) contains \(\mathbb Q(\mu_8)\)
and \([K:\mathbb Q]=d=|H|\). Remove every ramified rational prime
in the common set \(S\), and use the primitive characters on the
remaining Euler factors. If the Frobenius element at \(p\notin S\)
has order \(f\), the elementary character identity is
\[
 \prod_{\xi\in H}(1-\xi(p)T)=(1-T^f)^{d/f}.
\]
It proves, first for \(\Re u>1\) and then meromorphically,
\[
 \zeta_K^S(u)=\prod_{\xi\in H}L_{\mathbb Q}^S(u,\xi).
\]
Define the \emph{plain companion product}
\(C_\chi(u)=\prod_{\xi\in H\setminus\{\chi\}}
L_{\mathbb Q}^S(u,\xi)\). Equation
\eqref{eq:sp4-principal-inverse} gives the exact identity
\begin{equation}\label{eq:sp4-character-extraction}
 C_\chi(u)E_\psi((u-2)/4)
   =\frac{B((u-2)/4)}{L_{\mathbb Q}^S(u,\chi)}.
\end{equation}
This is a meromorphic identity, including at shared zeros: their orders
cancel as \(\sum_{\xi\in H}m_\xi-
\sum_{\xi\ne\chi}m_\xi=m_\chi\). No nonvanishing assumption on
the companion factors is made. The existing maintained
\texttt{CarmichaelDirichletProduct.product\_order\_formula}
and \texttt{CarmichaelWeightedSignal.weighted\_limit} implement
the product-order and signed local-readout principles; the new
adelic realization is not asserted by those declarations.

Here is the complete physical inverse consumer. Let \(a_K(n)\) be
the sum of \(\mu_K(\mathfrak a)\) over ideals of norm \(n\) prime to
\(S\), and let \(\xi_1,\ldots,\xi_{d-1}\) enumerate
\(H\setminus\{\chi\}\). For \(w\in C_c^\infty(0,\infty)\),
\begin{align*}
 &\sum_{a,b_1,\ldots,b_{d-1}\ge1}
   a_K(a)\prod_{j=1}^{d-1}\xi_j(b_j)
   \mathbf1_{(ab_1\cdots b_{d-1},S)=1}
   w(ab_1\cdots b_{d-1}/X)\\
 &\hspace{30mm}=\sum_{\substack{n\ge1\\(n,S)=1}}
      \mu(n)\chi(n)w(n/X).
\end{align*}
Each sum is finite because of the displayed common support.
This follows by the Euler-product coefficient identity
\(a_K*\xi_1*\cdots*\xi_{d-1}=\mu\chi\), with zero extension at
\(S\). No pairwise coprimality or independent factor cutoff is added.
All unit rows and all convolution cross terms remain present.
For \(K=\mathbb Q(\mu_8)\) and \(\chi=1\), this uses just the three
fixed quadratic companions. At \(p\equiv1\pmod8\), the local
calculation is \((1-T)^4(1-T)^{-3}=1-T\). At any other odd
residue it is
\((1-T^2)^2/((1-T)(1+T)^2)=1-T\).
Thus the cancellation is coefficientwise at every prime power,
not an absolute-value estimate on the separate factors.

Differentiation gives a second full consumer, with \(u=4s+2\):
\begin{equation}\label{eq:sp4-signed-log-extraction}
 \zeta_K^S(u)E_\psi'(s)-B'(s)
 +4B(s)\sum_{\xi\in H\setminus\{\chi\}}
       \frac{(L_{\mathbb Q}^S)'(u,\xi)}{L_{\mathbb Q}^S(u,\xi)}
 =-4B(s)\frac{(L_{\mathbb Q}^S)'(u,\chi)}
                    {L_{\mathbb Q}^S(u,\chi)}.
\end{equation}
The derivative \(B'\), the chain factor \(4\), every plain companion
and every higher prime power are compulsory. For nonprincipal \(\chi\),
the principal companion cancels the native principal term; for principal
\(\chi\) that term survives and must be subtracted in the physical sum.
There is no division by \(B\) away from a chosen nonvanishing window.

\paragraph{The restricted cubic character example has explicit costs.}
For a primitive cubic character of prime conductor \(q\), take
\(H=\{\psi\chi^j:\psi\bmod8,\ j=0,1,2\}\). Then
\(d=12\), the conductor product is \(2^{24}q^8\), and the
archimedean degree is \(12\). The native logarithmic derivative has
first-prime coefficient
\[
 12\log p\,\mathbf1_{p\equiv1\ (8),\ \chi(p)=1},\qquad p\notin S.
\]
The conductor product follows from the modulo-\(8\) conductors
\(1,4,8,8\), and the two nontrivial cubic powers of conductor \(q\).
Its height-dependent analytic conductor has order
\(2^{24}q^8(1+|t|)^{12}\), up to the fixed convention for gamma factors.
Extracting just \(\chi\) retains eleven companions. Estimating them
before exact coefficient cancellation charges every multiplied length,
not only one factor. Restriction to this prime class is therefore an
explicit family adapter, not a free all-character estimate.

\paragraph{The full affine window transfer.}
The physical scale must be changed together with the detector.
In general suppose \(u=as+b\), with real \(a>0,b\), and an exact
meromorphic identity \(C(u)E(s)=B(s)/L(u)\).
For \(c>1\), set \(c'=(c-b)/a\). On finite vertical segments,
ordinary substitution gives
\begin{align*}
 P(X;T)&=\frac1{2\pi i}\int_{c-iT}^{c+iT}
    \frac{\widehat w(u)B((u-b)/a)}{L(u)}X^u\,du\\
 &=aX^b\frac1{2\pi i}
   \int_{c'-iT/a}^{c'+iT/a}
      E(s)C(as+b)\widehat w(as+b)(X^a)^s\,ds.
\end{align*}
Every factor, both endpoints and the Jacobian are retained.
The substitution is valid without an infinite-contour convergence
claim. Passing \(T\to\infty\), or shifting either contour, additionally
requires bounds for the complete integrand and the corresponding
horizontal and truncation remainders.

Write \(F(Y)\) for the full spectral integral on the second line,
when justified. Then \(P(X)=aX^bF(X^a)\).
An independent estimate \(F(Y)\ll Y^{r+\varepsilon}\) gives
\(P(X)\ll X^{b+ar+a\varepsilon}\). Under the detector's Mellin
continuation hypotheses, including the small-\(Y\) end, a zero
\(\rho\) with nonzero window multiplier is excluded whenever
\(\Re\rho>b+ar\). Thus the boundary budget is \(b+ar\), not
the numerical exponent imported from another observable.
If an estimate costs \(q^A\), uses a shifted length \(q^\lambda Y\),
and has normalization \(q^{-\nu}\), the complete budget for
\(q\asymp X^\delta\) is instead
\[
 b+ar+\delta(A+\lambda r-\nu).
\]
Such a moving-family conclusion requires simultaneous window and
signal bounds in precisely that range; fixed-\(q\) detection is
insufficient.

For \eqref{eq:sp4-character-extraction}, \(a=4,b=2\),
\(C=C_\chi\), and \(Y=X^4\). Hence the required spectral exponents
are
\[
 \begin{array}{c|ccc}
 \text{Euler boundary}&7/8&3/4&1/2\\\hline
 \text{spectral exponent }r&-9/32&-5/16&-3/8.
 \end{array}
\]
The integral includes \(B\), all companions and all section costs.
It is not the unweighted coefficient sum unless the multiplier is
removed with a separately justified window transfer. In particular,
the finite-point choice \(B(s_\rho)=1\) does not justify global
division by \(B\) or eliminate its vertical growth.

\paragraph{What exponent the alternative must actually prove.}
For each fixed \(\chi\), an independent bound
\[
 \sum_{\substack{n\ge1\\(n,S)=1}}\mu(n)\chi(n)w(n/X)
       \ll_{\chi,w,\varepsilon}X^{b+\varepsilon}
 \quad(X\ge2)
\]
for a set of smooth windows detecting every hypothetical zero implies
the boundary \(\Re\rho\le b\), by Mellin continuation of the reciprocal.
Equivalently one may bound the correctly centred logarithmic consumer.
The finite-point section argument removes a pointwise choice issue;
it does not give this bound. A uniform proof through Eisenstein series
still needs the physical Bruhat transformation, the continuous local
interface above, and all moving section and companion costs.
The usual Weyl reflection \(u\mapsto4-u\) supplies none of this saving.
The installed cubic estimate remains \(7/8\); \(b=3/4\) requires a
new signed arithmetic estimate for this consumer. No exponent is
imported from a different detector's normalization.
